\section{Empirical analysis}\label{sec:emp}

\fbox{TODO:}
In practice, we linearly interpolated and extrapolated $\watm$. 

\begin{table}
%\begin{wraptable}{r}{0.75\textwidth}
\caption{\label{tab:appli_stats}Daily statistics for January 2018.}
\centering
\begin{tabular}{lrrrr}
\toprule
 & Mean & Std & Min & Max\\
\midrule
Max moneyness & 1.274 & 0.017 & 1.239 & 1.304\\
Min moneyness & 0.186 & 0.019 & 0.174 & 0.241\\
Number of contracts & 3434.000 & 193.934 & 3063.000 & 3747.000\\
Number of maturities & 30.722 & 0.461 & 30.000 & 31.000\\
\bottomrule
\end{tabular}
%\end{wraptable}
\end{table}

In this section, we apply our approach to modeling implied volatility surfaces extracted from S\&P500 options prices between.
More specifically, we use a dataset containing the price of all options traded during January 2018, namely 61812 contracts, and \Cref{tab:appli_stats} displays daily summary statistics.

We perform the following exercise for each day in the sample.
First, we split the daily sample into a training and a testing set.
Second, we fit the model on the training set and evaluate its performance on the testing set.
As described in~\Cref{tab:intraextra}, we use two different configurations for training and testing.
In the interpolation setting, for each maturity, we randomly select half of the contracts.
As such, we also sample options that are far out or in the money for training, and the testing error represents the approximation error for the range of moneyness that are actually observed.
In the extrapolation setting, for each maturity, we select 70\% of the contract that have an absolute Call option Delta in the range $[0.02,0.98]$, see the Appendix~B.2.
This second filter is equivalent to selecting approximately 50\% of all the contracts but with more observation around the money.
Thus, we do not select options that are far out or in the money for training, and the testing error measures how well our model extrapolates.
Finally, we again use the three set of conditions described in~\Cref{tab:conditions} to study how the arbitrage-related penalties affect the results.

\begin{table}
%\begin{wraptable}{r}{0.6\textwidth}
\caption{Settings for the training and testing sets.\label{tab:intraextra}}
\centering
\begin{tabular}{ l c c   }
  \hline
    & Call Delta & Train/(Train + Test) \\
\hline
   Interpolation & 0--1 & 0.5 \\
\hline
   Extrapolation & 0.02--0.98 & $\approx$ 0.5 \\
\hline
\end{tabular}
%\end{wraptable}
\end{table}

In~\Cref{tab:appli_loss}, we present our results.
First, we describe the RMSE and MAPE.
As expected, training errors are generally below testing errors.
Somewhat surprisingly, extrapolation errors are most of the time smaller than interpolation errors.
This may be because the surface is highly non-linear around the money, i.e.\ around $k=0$, while the surface wings tend to be more linear.
While MAPE errors might seem large, it should be noted that, for some contracts, the implied volatility is extremely small and thus even tiny deviations significantly impact this error measure.
Interestingly, we see that the second condition (i.e., strong enforcement of the no-arbitrage constraints) also leads to the smallest interpolation and extrapolation errors.
Furthermore, comparing the first and third conditions, we see that even a mild no-arbitrage constraint provides improvements over simply trying to minimize the RMSE and MAPE.
Regarding CAL and BUT, we see that the models resulting from condition 1 and 2 are essentially arbitrage-free.
As for the model resulting from no enforcement of the constraints (i.e., condition 3), it provides calendar-spread arbitrage opportunities both in interpolation and extrapolation, as well as butterfly arbitrage opportunities in extrapolation.

Time-series of the daily losses can be found in Appendix~C.3.
\vspace{-0.2cm}

\begin{table}[!h]

\caption{Mean loss (with standard deviation).
    All numbers are multiplied by 100.
    \label{tab:appli_loss}}
\centering
\resizebox{\linewidth}{!}{
\begin{tabular}{llllllll}
\toprule
\multicolumn{2}{c}{ } & \multicolumn{2}{c}{Condition 1} & \multicolumn{2}{c}{Condition 2} & \multicolumn{2}{c}{Condition 3} \\
\cmidrule(l{3pt}r{3pt}){3-4} \cmidrule(l{3pt}r{3pt}){5-6} \cmidrule(l{3pt}r{3pt}){7-8}
 &  & Interpolation & Extrapolation & Interpolation & Extrapolation & Interpolation & Extrapolation\\
\midrule
 & train & 3.12 (3.04) & 4.79 (5.82) & 2.47 (2.98) & 1.56 (3.58) & 5.46 (1.91) & 8.1 (5.7)\\
\cmidrule{2-8}
\multirow[t]{-2}{*}{\raggedright\arraybackslash RMSE} & test & 9.48 (7.52) & 4.87 (5.91) & 7.66 (6.95) & 1.55 (3.47) & 14.27 (4.67) & 8.2 (5.75)\\
\cmidrule{1-8}
 & train & 16.04 (15.67) & 16.39 (19.88) & 12.81 (15.46) & 5.42 (12.43) & 27.96 (10.09) & 27.66 (19.43)\\
\cmidrule{2-8}
\multirow[t]{-2}{*}{\raggedright\arraybackslash MAPE} & test & 26.4 (22.65) & 16.66 (20.11) & 21.61 (21.91) & 5.48 (12.26) & 42.14 (13.54) & 27.76 (19.45)\\
\cmidrule{1-8}
 & train & 0.01 (0.02) & 0.02 (0.03) & 0.03 (0.07) & 0.07 (0.18) & 4.24 (2.75) & 2.39 (1.78)\\
\cmidrule{2-8}
\multirow[t]{-2}{*}{\raggedright\arraybackslash CAL} & test & 0.01 (0.02) & 0.02 (0.03) & 0.03 (0.07) & 0.07 (0.18) & 4.24 (2.75) & 2.39 (1.78)\\
\cmidrule{1-8}
 & train & 0 (0) & 0 (0) & 0 (0) & 0 (0) & 0.86 (2.06) & 3.36 (14.06)\\
\cmidrule{2-8}
\multirow[t]{-2}{*}{\raggedright\arraybackslash BUT} & test & 0 (0) & 0 (0) & 0 (0) & 0 (0) & 0.86 (2.06) & 3.36 (14.06)\\
\bottomrule
\end{tabular}}
\end{table}



